    % !TeX root = ../main.tex
\section{Выводы}
В ходе выполнения работы получены аналитические выражения и построены графики зависимостей СКО фильтрации доплеровской частоты $\sigma_\Omega$ и эквивалентной шумовой полосы $\Delta F_{\text{ЧАП}}$ от отношения сигнал/шум $q_{c/n_0}$ и среднеквадратического ускорения $\sigma_a$. Показано, что при улучшении помеховой обстановки или росте среднеквадратического ускорения птимальный фильтр расширяет полосу пропускания.

Для задачи слежения в дискретном времени записаны рекуррентные уравнения оптимального фильтра Калмана второго порядка и проведено статистическое моделирование работы системы ЧАП на интервале $1000\text{ с}$. Время вхождения в синхронизм по дисперсии ошибки составило около $0.5\text{ с}$.

Сравнение теоретического значения $D_{11} = 132.26\text{ (рад/с)}^2$ и выборочной дисперсии ошибки фильтрации $\hat{D}_{11} = 133.08\text{ (рад/с)}^2$ в установившемся режиме показало их высокое согласование, что подтверждает оптимальность синтезированного алгоритма.
